\subsection{Criteria for a Quaternion Algebra to Be a Field}

Let \(\alpha,\beta,\gamma\) be elements of the field K, and let F be the quaternion algebra of type \((\alpha,\beta,\gamma)\) . As in No.~1, we denote the canonical basis of F by \((1,i,j,k)\) and the subalgebra \(K+Ki\) of F by E.

PROPOSITION 5. --- The following properties are equivalent:

\begin{enumerate}
\def\labelenumi{(\roman{enumi})}
\item
  The quaternion algebra \(\mathbf{F}\) is a field.
\item
  There is no nonzero element \(q \in \mathrm{F}\) such that \(\mathrm{T}_{\mathrm{F}}(q) = \mathrm{N}_{\mathrm{F}}(q) = 0\) .
\item
  Every element of F of square zero is zero.
\item
  There does not exist any nonzero vector \((x,y,z,t)\) in \(K^{4}\) such that
\end{enumerate}

\[
x ^ {2} + \beta x y - \alpha y ^ {2} - \gamma (z ^ {2} + \beta z t - \alpha t ^ {2}) = 0.
\]

\begin{enumerate}
\def\labelenumi{(\alph{enumi})}
\setcounter{enumi}{21}
\tightlist
\item
  There does not exist any nonzero vector \((x,y,z)\) in \(K^{3}\) such that
\end{enumerate}

\[
x ^ {2} + \beta x y - \alpha y ^ {2} - \gamma z ^ {2} = 0.
\]

\begin{enumerate}
\def\labelenumi{(\roman{enumi})}
\setcounter{enumi}{5}
\tightlist
\item
  The quadratic algebra E is a field, and \(\gamma\) is not the norm of an element of E.
\end{enumerate}

An element q of F is invertible if and only if \(\mathrm{N}_{\mathrm{F}}(q)\) is not 0. The equivalence of (i) and (iv) therefore follows from formula (31) of III, §2, No.~5, p.~445; it is clear that (i) implies (iii) and that (iv) implies (v).

The equivalence of (ii) and (iii) follows from Remark 3 of VIII, p.~362.

Suppose that F is not a field. If \(\gamma(4\alpha+\beta^{2})\neq0\) , then the algebra F is central simple of degree 4 over K; it is isomorphic to the algebra \(\mathbf{M}_{2}(\mathrm{K})\) (VIII, p.~120, Theorem 1) and therefore contains a nonzero element of square zero. If \(\gamma=0\) , then we have \(j^{2}=0\) . If \(4\alpha+\beta^{2}=0\) , then we have \((2i-\beta)^{2}=0\) and \(2i-\beta\neq0\) if K has characteristic different from 2. Finally, if K has characteristic 2 and \(\beta\) is zero, then we have \(\mathrm{T}_{\mathrm{F}}(q)=0\) for every \(q\in F\) (III, §2, No.~5, p.~445, formula (31)). Since F is not a field, there exists a nonzero element q of F such that \(\mathrm{N}_{\mathrm{F}}(q)=0\) ; we have \(q^{2}=0\) . This proves the implication \((\mathrm{iii})\Rightarrow(\mathrm{i})\) .

Let \(q = x + yi\) be an element of E. We have \(\mathrm{N}_{\mathrm{E / K}}(q) = x^2 + \beta xy - \alpha y^2\) . Suppose that property (v) holds. We have \(\mathrm{N}_{\mathrm{E / K}}(q) - \gamma \neq 0\) and \(\mathrm{N}_{\mathrm{E / K}}(q) \neq 0\) if \(q \neq 0\) , and therefore (vi).

Finally, suppose that property (vi) holds. Let q be a nonzero element of F; we write it as \(u + vj\) with u and v in E. If v is zero, then q is invertible. If v is not zero, then we have \(\mathrm{N}_{\mathrm{F}}(q) = \mathrm{N}_{\mathrm{F}}(v)\mathrm{N}_{\mathrm{F}}(v^{-1}u + j) = \mathrm{N}_{\mathrm{E}/\mathrm{K}}(v)\big(\mathrm{N}_{\mathrm{E}/\mathrm{K}}(v^{-1}u) - \gamma\big)\) by III, §2, No.~5, p.~443, formula (24). Since \(\gamma\) is not a norm, \(\mathrm{N}_{\mathrm{F}}(q)\) is not zero, and q is invertible.

Remark. --- Suppose that the quaternion algebra F is a field. It follows from the equality \(j^{2} = \gamma\) that we have \(\gamma \neq 0\) . By Proposition 2 of VIII, p.~363, the center of F is equal to K unless K has characteristic 2 and \(\beta\) is zero, in which case the algebra F is commutative.

